\section{Response Geometry and Command Information in \texorpdfstring{$\pi_0$}{the Second Policy}}
\label{sec:supp-pi0-internals}

Do different perturbations also change the features of $\pi_0$ along shared directions? Yes. On held-out tasks and families, 16 shared directions retain 70--77\% of the response energy at L9, 39--50\% at L1 and 43--52\% at L18, against 1--28\% for clean-feature and about 1.5\% for random directions, in every family (Table~\ref{tab:pi0-retention}). This supports claim C1 of the main paper (Sec.~\ref*{M-sec:shared}), which states that the result holds in $\pi_0$. The command readout qualifies it: in $\pi_0$ the shared coordinates carry command information at L18 but predict commands no better than clean-feature coordinates.

\paragraph{What was compared.}
We recorded $\pi_0$ at the first query on its paired LIBERO-Plus inputs and compared shared, clean-feature and random directions, 16 each. In Table~\ref{tab:pi0-retention} each row fixes the layer and family; higher retention means a better reconstruction of the same centered response (task-equal means over 28 task components, 2{,}000 whole-task bootstrap draws).

\paragraph{The shared directions carry command information at L18 but do not beat clean-feature directions.}
At L18, a linear readout of the shared coordinates predicts the change in the first five commands (translation) better than the magnitude baseline and better than random coordinates, with intervals above zero in 4 of 4 families for each comparison; the corresponding counts are 0 of 4 at L1 and 0 of 4 at L9. Against the full hidden state or the centered clean-feature coordinates, no interval is above zero at any of the three layers (0 of 4 each), and adding the shared coordinates to first flow gives no interval above zero. The readout thus finds command information in the shared directions at L18 but does not favor them over clean-feature coordinates.

\paragraph{Recording $\pi_0$ at the same computation stage.}
The $\pi_0$ plan has 50 actions, and its action expert prepends one state token; the recorded state is $[10,3,1,51,1024]$, and the analysis uses the 50 plan positions. The original $\pi_0$ paired executions stored no internal states, so the policy was rerun on their recorded first-query inputs and policy noise, on different hardware. Regenerated raw actions differ from the recorded actions by at most 0.006--0.008 in absolute value. Input reconstruction is deterministic, and two independent reruns of the same run on the same hardware agree exactly.

\paragraph{Fitting the directions and the command readouts.}
We use the joint task/family split and affine readout of Sec.~\ref{sec:current-basis-readout}. For each held-out family, five task-component folds exclude both that family and the test fold from fitting; testing uses their intersection. Tasks that share a clean scene stay in one component, and common tasks keep the $\pi_{0.5}$ fold assignments. Each split fits position means and shared and clean-feature directions anew on its admitted $\pi_0$ fitting recordings. Fitting weights tasks equally, then families, settings, noise draws and positions within their parents.

At the first generation step, removing the state token leaves 50 ordered positions with 1{,}024 channels. Each set of 16 directions shares its channel basis across positions, giving 800 coefficients per layer; positionwise magnitudes, full hidden responses and first-flow differences use 50, 51{,}200 and 1{,}600 features. First flow plus the shared coordinates uses 2{,}400 features. Each block is centered on its weighted fitting mean and divided by one fitting RMS scale; affine ridge uses an unpenalized intercept and fixed $\lambda=0.01$. The combined readout scales its two blocks separately. Layers are fitted separately, with position coordinates kept in each feature vector. We score the regenerated first five commands (after the output adapter) separately for translation, rotation and continuous gripper.

\paragraph{Admission rule for rerun recordings.}
A rerun was admitted if its actions stayed within twice the largest difference observed when the original $\pi_0$ runs were repeated on the same input across machines: raw-action maximum absolute difference $\le 0.018$; first-five translation RMS $\le 3.1\times10^{-3}$, rotation RMS $\le 2.1\times10^{-3}$, gripper maximum $\le 3.9\times10^{-3}$; sign agreement where the recorded value exceeds three times the threshold; no NaN. The rule depends on no outcome. It replaced an exact-match criterion after the discrepancy was observed.

\paragraph{Admitted pairs and exclusions.}
Of 808 runs, 769 were admitted and 39 excluded without retry: 7 runs of four Light settings were missing for technical reasons and 32 exceeded the tolerance (25 on the gripper threshold, 7 on the raw-action threshold, largest 0.27). This left 370 of the 404 planned clean/perturbed pairs on 28 tasks. One further pair was dropped, not retried, because the recorded state of an admitted run failed to write, so 369 pairs enter the analysis. Background, Camera, Light and Sensor Noise contribute 91/89/83/106 pairs on 28/28/27/28 task components. The exclusions spread over 21 family--task cells and all four families (Background 6 of 192 runs, Camera 15 of 204, Light 14 of 192 including the 7 technical, Sensor Noise 4 of 220); 7 of the 31 affected settings lost more than one run. Every statistic compares clean and perturbed recordings from the same rerun; the original actions serve only as the reference for admission, and admitted reruns lie within the original cross-machine scale rather than matching exactly.

\begin{table*}[t]\centering\small
\caption{\textbf{In $\pi_0$, the 16 shared directions retain far more of an unseen response than clean-feature or random directions.} Share of the response energy around the fitting-data mean (\%) at the first query, task-equal mean on each family's measured pairs with pointwise 95\% whole-task bootstrap intervals; 369 pairs on 28 task components overall (family counts in text); directions fitted without the test tasks and family. The recordings come from a rerun of $\pi_0$ on different hardware, admitted by the action-difference criterion in the text.}
\label{tab:pi0-retention}
\begin{tabular}{llrrr}\toprule
Layer & Family & Shared & Clean-feature & Random\\\midrule
L1 & Background & 44.4 [41.1, 48.0] & 2.0 [1.8, 2.1] & 1.4 [1.4, 1.5]\\
L1 & Camera & 49.2 [43.9, 54.8] & 2.2 [2.1, 2.3] & 1.5 [1.4, 1.6]\\
L1 & Light & 38.5 [35.9, 41.4] & 2.2 [2.1, 2.2] & 1.5 [1.5, 1.6]\\
L1 & Sensor Noise & 50.1 [46.2, 54.0] & 2.0 [1.8, 2.2] & 1.4 [1.3, 1.4]\\
L9 & Background & 72.2 [69.4, 75.1] & 1.0 [0.9, 1.1] & 1.8 [1.7, 1.9]\\
L9 & Camera & 77.2 [74.1, 80.0] & 1.0 [0.9, 1.1] & 1.7 [1.6, 1.8]\\
L9 & Light & 71.9 [69.6, 74.1] & 1.1 [1.0, 1.2] & 1.9 [1.8, 2.0]\\
L9 & Sensor Noise & 70.3 [66.3, 74.0] & 1.7 [1.6, 1.8] & 1.9 [1.8, 2.0]\\
L18 & Background & 43.7 [39.9, 48.1] & 21.5 [18.8, 24.4] & 1.5 [1.4, 1.6]\\
L18 & Camera & 51.9 [47.8, 55.5] & 28.1 [24.0, 32.0] & 1.5 [1.4, 1.6]\\
L18 & Light & 49.2 [46.2, 52.5] & 17.5 [15.3, 20.0] & 1.4 [1.3, 1.5]\\
L18 & Sensor Noise & 43.0 [39.2, 47.1] & 21.5 [17.7, 25.9] & 1.5 [1.4, 1.6]\\\bottomrule
\end{tabular}
\end{table*}
